\chapter{La première épreuve}
\label{chap:premiere-epreuve}

Les années passèrent et Khaer Ghal grandit comme grandissaient les autres
enfants du village, au rythme des tâches, des entraînements et des
enseignements que les adultes considéraient indispensables à leur survie.

Très tôt, il apprit à reconnaître les traces laissées par les animaux autour
des habitations, à distinguer les plantes que l'on pouvait cueillir de celles
qu'il valait mieux éviter et à retrouver son chemin lorsqu'il s'éloignait dans
les bois. On lui enseigna à courir sans épuiser inutilement ses forces, à
grimper, à tomber sans se blesser et à supporter la faim, le froid ou la
fatigue sans s'en plaindre. Toutes ces connaissances faisaient partie de la
vie quotidienne du village et personne ne songeait réellement à les séparer
les unes des autres : savoir survivre était aussi important que savoir se
battre.

Le combat occupait néanmoins une place particulière dans son éducation. Les
premières leçons se faisaient à mains nues, puis avec des bâtons de bois. Les
enfants apprenaient à frapper, à esquiver, à conserver leur équilibre et,
surtout, à continuer lorsqu'un premier coup les envoyait au sol. Plus tard
venaient les armes. Khaer Ghal était encore trop jeune pour manier celles des
adultes, mais on lui enseigna l'usage d'une lame courte ainsi que les rudiments
du combat contre les animaux et les créatures qui vivaient autour du village.

Son père veillait à ce que le fils du chef ne bénéficie d'aucune indulgence.
Il se montrait même parfois plus exigeant avec lui qu'avec les autres enfants.
Khaer Ghal portait son nom et, à ses yeux, cela impliquait qu'il fasse honneur
à la place occupée par leur famille au sein du clan. Là où l'erreur d'un autre
enfant pouvait être simplement corrigée, la sienne était plus difficilement
acceptée. S'il devait un jour prendre sa place parmi les guerriers, il devrait
la mériter comme les autres ; et puisqu'il était le fils du chef, son père
attendait même de lui qu'il se montre capable d'en faire davantage.

Ce fut ainsi qu'à seulement huit ans Khaer Ghal fut autorisé à se présenter à
sa première véritable épreuve.

C'était tôt. La plupart des enfants disposaient encore de quelques années
d'entraînement avant d'être confrontés seuls à une créature, mais son père
avait estimé qu'il était prêt. Khaer Ghal ignorait s'il partageait réellement
cette certitude ; il savait seulement qu'il n'avait aucune intention de donner
à son père une raison d'en douter.

La tradition voulait qu'un jeune Orc prouve un jour qu'il était capable
d'appliquer seul ce qu'on lui enseignait depuis son enfance. Il ne s'agissait
pas encore d'affronter un autre guerrier, mais une créature capturée dans les
environs du village et choisie de manière à représenter un danger réel sans
être considérée comme hors de portée d'un enfant correctement entraîné.

Ce matin-là, une bonne partie du village s'était rassemblée autour de la grande
fosse creusée au centre de la colline. Khaer Ghal connaissait parfaitement
l'endroit. Il avait passé des années à observer les combats depuis les
aménagements de bois et de pierre disposés le long du rebord, plusieurs mètres
au-dessus du sol de l'arène. Vu d'en haut, entouré des autres habitants, tout
lui avait toujours semblé relativement simple.

Depuis le fond de la fosse, les choses paraissaient différentes.

Un guerrier qui avait participé à son entraînement l'accompagna jusqu'à l'un
des passages aménagés dans la paroi et lui tendit une dague adaptée à sa
taille. Khaer Ghal referma les doigts autour de la poignée, descendit les
dernières marches taillées dans la roche puis s'avança sur la terre battue.
Derrière lui, une lourde grille referma l'accès.

À l'autre extrémité de la fosse, un second passage s'ouvrit.

La créature hésita quelques secondes avant de quitter l'ombre. Elle avançait
sur quatre pattes, le corps bas, souple et nerveux, avec une longue queue qui
accompagnait chacun de ses mouvements. Elle n'était ni particulièrement grande
ni impressionnante aux yeux des guerriers qui la regardaient depuis le haut,
mais sa mâchoire large et les longues canines qui dépassaient légèrement de sa
gueule suffirent à rendre la dague de Khaer Ghal beaucoup moins rassurante.

À huit ans, elle lui semblait largement assez dangereuse.

Il raffermit sa prise et se força à reprendre ce qu'on lui avait répété pendant
des mois : ne pas attaquer le premier, observer, attendre que la créature
s'engage, éviter la gueule et profiter de l'ouverture lorsqu'elle se
présenterait. Pendant les entraînements, toutes ces règles avaient toujours
semblé parfaitement logiques. Elles devenaient beaucoup moins évidentes
lorsqu'une véritable créature tournait lentement autour de lui en cherchant le
meilleur moment pour bondir.

Khaer Ghal pivota avec elle, sans la quitter des yeux.

%\illustration
%  {0200.png}
%  {L'épreuve de Khaer Ghal}
%  {epreuve}

Lorsqu'elle attaqua enfin, il reconnut suffisamment tôt le mouvement de ses
pattes arrière. Il se décala comme on le lui avait appris et frappa au passage.
La lame atteignit le flanc de l'animal, mais beaucoup moins profondément qu'il
ne l'avait espéré. La créature poussa un cri, pivota aussitôt et l'obligea à
reculer précipitamment.

Pendant une seconde, Khaer Ghal sentit pourtant monter en lui une satisfaction
presque triomphante. Il avait esquivé et il avait touché ; tout ce qu'on lui
avait enseigné fonctionnait donc réellement.

Le sourire qui commençait à se dessiner sur son visage disparut lorsque la
créature attaqua une seconde fois. Khaer Ghal voulut reproduire exactement le
même mouvement, mais cette fois la bête modifia sa trajectoire au dernier
instant. Elle le heurta de plein fouet et l'envoya rouler dans la poussière.
Sa dague lui échappa des doigts et s'immobilisa quelques pas plus loin, la
créature se retrouvant désormais entre lui et son arme.

Autour de la fosse, les cris et les encouragements changèrent brusquement de
ton. Khaer Ghal se releva aussi vite qu'il le put tandis que l'animal repartait
déjà vers lui. Il attendit jusqu'au dernier instant pour se jeter sur le côté ;
une patte accrocha son vêtement et le fit basculer une nouvelle fois, mais il
était désormais plus près de la dague.

Il tendit la main vers son arme et ses doigts en effleurèrent la poignée au
moment où une douleur brutale lui traversa tout le côté gauche du crâne.

La mâchoire de la créature venait de se refermer sur son oreille.

Khaer Ghal hurla et frappa instinctivement de son bras libre, sans parvenir à
faire lâcher prise à l'animal. La bête secoua violemment la tête et la douleur
devint fulgurante. Il sentit le cartilage céder sous les crocs et une partie
du bord supérieur de son oreille se déchirer. Toujours plaqué au sol, il
chercha désespérément sa dague à tâtons, referma enfin les doigts sur sa
poignée et frappa aussi fort qu'il le put.

La créature lâcha prise.

Khaer Ghal roula immédiatement hors de sa portée et porta une main à son
oreille gauche. Lorsqu'il la retira, sa paume était couverte de sang. La
morsure n'avait pas seulement ouvert la peau : le bord supérieur de l'oreille
était profondément entaillé et un petit morceau avait été arraché, laissant
une encoche irrégulière dans le cartilage. Une douleur sourde pulsait jusque
dans sa mâchoire, mais il n'eut pas le temps de s'y attarder. En face de lui,
la créature était toujours debout.

Elle aussi saignait.

Cette fois, Khaer Ghal ne souriait plus.

Il se remit péniblement sur ses pieds, reprit sa garde et comprit quelque
chose que les entraînements ne lui avaient jamais réellement appris. Une
créature blessée n'abandonnait pas nécessairement. La douleur pouvait même la
rendre plus imprévisible, plus nerveuse et plus dangereuse.

Il cessa alors de chercher immédiatement à frapper. La bête s'approcha et
Khaer Ghal recula en observant chacun de ses mouvements. Lorsqu'elle bondit de
nouveau, il esquiva sans attaquer ; elle revint presque aussitôt et il recula
encore, davantage occupé à comprendre qu'à gagner du terrain.

Peu à peu, quelque chose se dessina dans les mouvements de l'animal. Avant
chaque bond, son corps se tassait légèrement, son poids passait sur les pattes
arrière et sa tête descendait une fraction de seconde. La première fois,
Khaer Ghal avait réagi trop tard pour s'en rendre compte ; désormais, il
commençait à attendre ce signal.

À l'attaque suivante, il se décala au moment où la créature prenait son appui,
mais n'alla pas assez vite. Une patte le heurta à l'épaule et le fit tourner
sur lui-même ; cette fois, pourtant, il parvint à rester debout. Emportée par
son élan, la bête passa devant lui et lui offrit pendant un instant son flanc.

Khaer Ghal frappa.

La lame entra plus profondément que la première fois. Dès que la créature se
retourna, il recula hors de portée plutôt que de chercher à profiter
immédiatement de son avantage.

Il commençait à comprendre.

Le combat se poursuivit encore plusieurs minutes. Khaer Ghal esquiva trop tôt
une attaque et reçut un nouveau coup de patte ; il tenta une riposte trop
précipitée et ne toucha que l'air. À plusieurs reprises, il dut simplement
mettre de la distance entre eux pour reprendre son souffle. Le sang continuait
de couler le long de son oreille et jusque sur son cou, son épaule le faisait
souffrir et la dague lui semblait de plus en plus lourde dans la main.

La créature, cependant, ralentissait elle aussi.

Lorsqu'elle se ramassa une nouvelle fois sur ses pattes arrière, Khaer Ghal
reconnut immédiatement la séquence qu'il avait observée pendant tout le
combat : le déplacement du poids, la tête qui descendait, puis l'instant où
tout son corps se préparait à partir vers l'avant.

Cette fois, il attendit.

La créature bondit.

Khaer Ghal ne bougea qu'au dernier moment. Il se décala juste assez pour éviter
la mâchoire et, au lieu de continuer à reculer, se jeta contre son flanc en
utilisant l'élan de l'animal contre lui.

Sa dague s'enfonça profondément.

La créature poursuivit encore sa course sur quelques pas avant de vaciller,
puis de s'effondrer dans la poussière.

Khaer Ghal resta immobile, les deux mains serrées autour de la poignée de son
arme, incapable pendant quelques secondes de détourner les yeux du corps
étendu devant lui. Les cris du village lui parvenaient comme s'ils venaient de
très loin ; il n'entendait réellement que sa propre respiration et les
battements de son cœur.

Puis le bruit de ceux qui observaient depuis le haut de la fosse lui revint
brusquement.

Il avait gagné.

\scenebreak

Le même guerrier qui lui avait remis la dague descendit dans l'arène lorsque
la grille fut de nouveau ouverte. Khaer Ghal s'attendait presque à l'entendre
énumérer immédiatement toutes les erreurs commises pendant le combat, car il
en avait déjà identifié suffisamment lui-même pour savoir que la liste serait
longue.

Le guerrier passa pourtant devant lui sans commentaire et s'accroupit près de
la créature. Après avoir sorti son couteau, il travailla quelques instants
autour de la mâchoire jusqu'à pouvoir en extraire l'une des longues canines
qui avaient tant impressionné Khaer Ghal au début du combat. Il l'essuya
sommairement avant de revenir vers lui.

Son regard se posa alors sur l'oreille gauche du garçon.

La morsure avait laissé des traces qu'il conserverait probablement toute sa
vie. Le cartilage était entaillé à plusieurs endroits et une petite partie du
bord supérieur de l'oreille avait été arrachée, donnant désormais à son
contour une forme irrégulière. Le guerrier examina la blessure quelques
instants, puis regarda la longue dent qu'il tenait encore dans sa main.

Selon une ancienne tradition du village, un enfant devait conserver sur lui
une partie de la première créature qu'il avait vaincue seul. Selon l'animal,
il pouvait s'agir d'une griffe, d'un morceau de corne ou d'une dent, autant de
trophées destinés à rappeler l'épreuve par laquelle le jeune Orc avait gagné
le droit de franchir une nouvelle étape de son apprentissage.

Dans le cas de Khaer Ghal, l'endroit où porter le sien semblait presque avoir
été choisi par la créature elle-même.

Il fallut néanmoins laisser à la blessure le temps de se refermer. Pendant
plusieurs jours, la canine resta donc de côté tandis que son oreille était
nettoyée et soignée. Lorsque le cartilage eut suffisamment cicatrisé, la dent
fut légèrement retaillée à sa base. Une petite encoche y fut creusée afin
qu'elle puisse rester solidement en place, puis elle fut passée directement
à travers l'oreille gauche de Khaer Ghal, à proximité des marques laissées par
la morsure.

Aucun métal, aucune parure ne venait l'accompagner. La canine était portée
telle qu'elle avait été prise sur la créature, à peine travaillée pour pouvoir
tenir en place. Sa teinte claire contrastait avec la peau gris-bleu de
Khaer Ghal et suivait en partie la forme de son oreille. Juste au-dessus,
l'irrégularité du cartilage et la portion manquante rappelaient immédiatement
d'où provenait ce trophée.

Khaer Ghal l'observa longuement la première fois qu'on lui permit de voir le
résultat. Son oreille ne retrouverait jamais sa forme d'origine, mais cela ne
semblait guère le préoccuper.

La dent, elle, lui plaisait beaucoup.

\scenebreak

En réussissant cette épreuve, Khaer Ghal avait également franchi une étape
importante aux yeux du village. Il n'était évidemment pas encore un guerrier,
mais il n'était plus considéré tout à fait comme l'un des plus jeunes enfants
du clan. Il pouvait désormais participer à des entraînements plus exigeants,
recevoir de nouvelles responsabilités et commencer à apprendre ce que l'on
réservait jusque-là aux plus âgés.

Pendant les jours qui suivirent, il prit cette nouvelle position très au
sérieux.

Peut-être même un peu trop.

Il parcourait le village avec une fierté qu'il ne faisait aucun effort pour
dissimuler et trouvait étonnamment souvent une raison de tourner la tête de
manière à présenter son profil gauche. La partie manquante de son oreille et
les cicatrices encore récentes auraient pu rappeler à n'importe qui la
violence de son épreuve ; pour Khaer Ghal, elles semblaient surtout mettre
davantage en valeur la longue canine qui la traversait désormais.

Il l'avait gagnée.

À huit ans, cela suffisait largement à justifier de la montrer à tout le
village.

Son père ne lui accorda guère plus qu'un signe approbateur lorsqu'il le vit
ainsi exhiber son trophée. Sa mère, elle, leva les yeux au ciel devant tant de
fierté, sans parvenir tout à fait à dissimuler son sourire.

Khaer Ghal s'en satisfit parfaitement.

Il avait réussi sa première épreuve.